\documentclass[14pt]{extarticle}
\usepackage{fontspec}
\usepackage{polyglossia}
\usepackage{amsmath}
\usepackage{amssymb}
\usepackage{xcolor}
\usepackage{fancyhdr}
\usepackage{graphicx}
\usepackage{listings}
\usepackage[most]{tcolorbox}
% \usepackage{setspace}
\usepackage{pifont}
\usepackage{enumitem}
\usepackage{titlesec}
\usepackage[bottom]{footmisc}
\usepackage{titling}
\usepackage{minted}
\usepackage{etoolbox}
\usepackage{array}
\usepackage{longtable} % page-breakable tables (must load before polyglossia/bidi RTL setup)
\usepackage{extsizes}
\usepackage{float}
\usepackage{booktabs}
\usepackage{hyperref}
\usepackage[onehalfspacing]{setspace}
\newif\ifwhiteblack
\ifdefined\setwhiteblack
  \setwhiteblack
\fi

\ifwhiteblack
  \usepackage[a4paper, landscape, top=1.25cm, bottom=1.25cm, left=1.25cm, right=1.25cm, headheight=1.5cm, headsep=0.5cm]{geometry}
\else
  \usepackage[a4paper, top=2.7cm, bottom=2.7cm, left=2.5cm, right=2.5cm, headheight=1.5cm, headsep=0.5cm]{geometry}
\fi
\usepackage{tikz}
\usepackage{datetime2}

\usetikzlibrary{automata, positioning, arrows.meta, arrows}
\providecommand{\LTcaptype}{table}

\newfontfamily\emoji{Segoe UI Emoji}

\pagestyle{fancy}

\setmainlanguage[numerals=western]{arabic}
\setotherlanguage{english}
\newfontfamily\arabicfont[Script=Arabic]{Calibri}
\newfontfamily\arabicfonttt[Script=Arabic]{Courier New}



\ifwhiteblack
\lstset{
  language=[Sharp]C,
  numbers=left,
  stepnumber=1,
  numbersep=8pt,
  frame=single,
  basicstyle=\ttfamily\small,
  keywordstyle=\bfseries,
  stringstyle=\ttfamily,
  commentstyle=\itshape
}
\else
\lstset{
  language=[Sharp]C,
  numbers=left,
  stepnumber=1,
  numbersep=8pt,
  frame=single,
  basicstyle=\ttfamily\small,
  keywordstyle=\color{blue},
  stringstyle=\color{red},
  commentstyle=\color{green!50!black}
}
\fi

\newif\ifdetailed
\ifdefined\setdetailed
  \setdetailed
\fi

\newif\ifwithsols
\ifdefined\setwithsols
  \setwithsols
\fi

\ifwhiteblack
  \AtBeginDocument{
    \LARGE
    \usemintedstyle{bw}
  }
\fi

% Default minted look: compact code that wraps long lines.
% Not applied in black-and-white landscape mode, which already uses \LARGE.
\ifwhiteblack\else
  \setminted{fontsize=\small, breaklines=true}
\fi

% unified tcolorboxes for articles
\ifwhiteblack
\tcbset{colback=white, colframe=black, fonttitle=\bfseries, boxrule=1pt}
\newtcolorbox{boxDef}[1][]{title={{\emoji📘} تعريف\ifx\\#1\\\else ~#1\fi :}}
\newtcolorbox{boxExercise}[1][]{title={{\emoji🧩} تمرين\ifx\\#1\\\else ~#1\fi :}}
\newtcolorbox{boxExample}[1][]{title={{\emoji📝} مثال\ifx\\#1\\\else ~#1\fi :}}
\newtcolorbox{boxNote}[1][]{title={{\emoji✨} ملاحظة\ifx\\#1\\\else ~#1\fi :}}
\newtcolorbox{boxAttention}[1][]{title={{\emoji🔔} تنبيه\ifx\\#1\\\else ~#1\fi :}}
\newtcolorbox{boxWarning}[1][]{title={{\emoji⚡} ملاحظة هامة\ifx\\#1\\\else ~#1\fi :}}
\newtcolorbox{boxSolution}[1][]{title={{\emoji✅} حل\ifx\\#1\\\else ~#1\fi :}}
\newtcolorbox{boxSymbol}[1][]{title={{\emoji🔣} رمز\ifx\\#1\\\else ~#1\fi :}}
\newtcolorbox{boxHint}[1][]{title={{\emoji💡} تلميح\ifx\\#1\\\else ~#1\fi :}}
\else
\tcbset{colback=white, colframe=black, fonttitle=\bfseries, boxrule=0.8pt}
\newtcolorbox{boxDef}[1][]{colback=blue!5!white,colframe=blue!75!black,
  title={{\emoji📘} تعريف\ifx\\#1\\\else ~#1\fi :}}
\newtcolorbox{boxExercise}[1][]{colback=cyan!5!white,colframe=cyan!70!black,
  title={{\emoji🧩} تمرين\ifx\\#1\\\else ~#1\fi :}}
\newtcolorbox{boxExample}[1][]{colback=yellow!5!white,colframe=orange!90!black,
  title={{\emoji📝} مثال\ifx\\#1\\\else ~#1\fi :}}
\newtcolorbox{boxNote}[1][]{colback=gray!10!white,colframe=black,
  title={{\emoji✨} ملاحظة\ifx\\#1\\\else ~#1\fi :}}
\newtcolorbox{boxAttention}[1][]{colback=magenta!10!white,colframe=magenta!80!black,
  title={{\emoji🔔} تنبيه\ifx\\#1\\\else ~#1\fi :}}
\newtcolorbox{boxWarning}[1][]{colback=red!5!white,colframe=red!75!black,
  title={{\emoji⚡} ملاحظة هامة\ifx\\#1\\\else ~#1\fi :}}
\newtcolorbox{boxSolution}[1][]{colback=green!5!white,colframe=green!60!black,
  title={{\emoji✅} حل\ifx\\#1\\\else ~#1\fi :}}
\newtcolorbox{boxSymbol}[1][]{colback=purple!5!white,colframe=purple!70!black,
  title={{\emoji🔣} رمز\ifx\\#1\\\else ~#1\fi :}}
\newtcolorbox{boxHint}[1][]{colback=teal!5!white,colframe=teal!60!black,
  title={{\emoji💡} تلميح\ifx\\#1\\\else ~#1\fi :}}
\fi


\ifwhiteblack
\tcbset{simplecode/.style={ colback=white, colframe=black, boxrule=1pt, arc=2pt, left=4pt,right=4pt,top=4pt,bottom=4pt}}
\else
\tcbset{simplecode/.style={ colback=gray!5, colframe=black!50, boxrule=0.4pt, arc=2pt, left=4pt,right=4pt,top=4pt,bottom=4pt}}
\fi
% Compact code boxes: tighter line spacing (the preamble uses onehalfspacing) and small vertical gaps.
% Must come AFTER the simplecode definition above.
\tcbset{simplecode/.append style={fontupper=\setstretch{0.95}, top=4pt, bottom=4pt, before skip=6pt, after skip=6pt}}
\newenvironment{boxCode}{\begin{tcolorbox}[simplecode]}{\end{tcolorbox}}

% Code cell used inside tables (keeps code left-to-right inside RTL tables).
% Use inside a fixed-width p{...} column: \begin{pycell}\begin{minted}...\end{minted}\end{pycell}
\newtcolorbox{pycell}{frame empty,colback=white,boxsep=0pt,left=2pt,right=2pt,top=2pt,bottom=2pt,width=\linewidth}

\newcolumntype{C}[1]{>{\centering\arraybackslash}p{#1}}

\makeatletter
\DTMnewdatestyle{dotted}{%
  \renewcommand{\DTMdisplaydate}[4]{\two@digits{##3}.\two@digits{##2}.##1}%
  \renewcommand{\DTMDisplaydate}{\DTMdisplaydate}%
}
\makeatother
\DTMsetdatestyle{dotted}

% redefine spaces after titles
\makeatletter
\renewcommand{\@maketitle}{%
  \begin{center}
    {\huge \bfseries \@title \par}%
    \vskip 0.2em % space between title and author
    {\large \@author \par}%
    % \vskip 0.2em % space between author and date
    % {\normalsize \@date \par}%
    {\small \hfill \begin{english}\DTMtoday\end{english} \par}%
  \end{center}
}
\makeatother

\fancyhf{} % clear default
\fancypagestyle{plain}{
  \fancyhf{}
  \fancyhead[L]{مدرسة التسامح الشاملة}
  % \fancyhead[L]{\includegraphics[height=1cm]{../../../images/logoTasamoh.png}}
  \fancyhead[R]{الأستاذ محمود اغبارية}
  \fancyfoot[C]{\thepage}
}

\fancyhead[L]{مدرسة التسامح الشاملة}
\fancyhead[R]{الأستاذ محمود اغبارية}
\fancyfoot[C]{\thepage}

\setcounter{tocdepth}{3} % only section subsection and subsubsection in TOC

\newcommand{\importpdfpage}[2]{%
    \thispagestyle{empty}
    \newgeometry{left=0cm,right=0cm,top=0cm,bottom=0cm}
    \includegraphics[page=#2,width=0.95\paperwidth,height=0.95\paperheight,keepaspectratio]{#1}
    \restoregeometry
}

\newcommand{\insertFullImg}[1]{%
    \noindent
    \makebox[\textwidth][c]{
        \includegraphics[width=0.9\paperwidth,keepaspectratio]{#1}%
    }%
}%



\title{كراس أسطول المجرّة — المرحلة الثانية: تطبيق الوراثة بين الفئات}

\begin{document}

\maketitle
\thispagestyle{fancy}

\begin{boxWarning}
\textbf{تعليمات هامة للعمل:}
تعتمد هذه المرحلة على الفئات التي بنيتموها في المرحلة السابقة (\textenglish{ScoutShip} و \textenglish{HibernationShip} و \textenglish{BasicShip}) — حدّثوا الفئات الموجودة ولا تبدأوا من الصفر. \\
في \textbf{كل تمرين}، يجب تحديث مقطع الفحص داخل العملية الرئيسية \textenglish{Main} وتشغيله للتأكد من صحة عمل الفئات والعمليات الجديدة التي كتبتموها.
\end{boxWarning}

% ==============================================================
\section*{المرحلة الثانية: تطبيق الوراثة بين الفئات}
% ==============================================================

\begin{enumerate}[label=\textbf{تمرين \arabic*:}, wide=0pt, leftmargin=*, itemsep=2.2em]

    \item \textbf{تطبيق الوراثة في \textenglish{BasicShip} و \textenglish{ScoutShip} و \textenglish{HibernationShip}} \\
    حدّثوا الفئة \textenglish{BasicShip} حسب قواعد البرمجة كائنية التوجه التي تعلمناها (صفات \textenglish{protected}، تغليف، عمليات بنائية)، بحيث تصبح \textbf{فئة أساسية (فئة أب)} لكل من الفئتين \textenglish{ScoutShip} و \textenglish{HibernationShip} اللتين سترثان منها. \\
    ابنوا من جديد الفئتين \textenglish{ScoutShip} و \textenglish{HibernationShip} بحسب قواعد البرمجة كائنية التوجه، بحيث \textbf{ترثان} من الفئة \textenglish{BasicShip} (استخدموا \textenglish{: base(...)} في العمليات البنائية، واحذفوا الصفات والعمليات المكررة). \\
    \textbf{الفحص في \textenglish{Main}:} \\
    تأكدوا أن كود الفحص في \textenglish{Main} يعمل بنجاح على الفئات الوارثة الجديدة.

    \item \textbf{إضافة فئة السفينة الأم (\textenglish{MotherShip})} \\
    أضيفوا فئة جديدة باسم \textenglish{MotherShip} تمثّل سفينة قيادة أم ترث من \textenglish{BasicShip} (أو من \textenglish{ScoutShip})، حسب قواعد البرمجة كائنية التوجه التي تعلمناها:

    {\centering
    \renewcommand{\arraystretch}{1.2}
    \small
    \begin{tabular}{|>{\raggedleft\arraybackslash}p{0.4\textwidth}|>{\raggedright\arraybackslash}p{0.5\textwidth}|}
    \hline
    \multicolumn{2}{|c|}{\textbf{الفئة الفرعية: \textenglish{MotherShip : BasicShip}}} \\
    \hline
    صفات موروثة من الفئة الأساسية: رقم الأسطول، رقم المحطة، رقم المركبة، رقم هوية القائد، وكمية الوقود. & \textenglish{(موروثة من BasicShip)} \\
    \hline
    صفات خاصة: اسم الأسطول، وحجم الاحتياطي الطارئ المسموح. & \textenglish{private string fleetName;} \newline \textenglish{private double emergencyReserve;} \\
    \hline
    عملية بنائية أولى تتلقى بيانات الأسطول والمحطة والمركبة، واسم الأسطول، ورقم هوية القائد. & \textenglish{public MotherShip(int fleetNumber, int baseNumber, int shipNumber, string fleetName, string captainID)} \\
    \hline
    عملية بنائية ثانية تتلقى بيانات الأسطول والمحطة والمركبة، ورقم هوية القائد، واسم الأسطول، وحجم الاحتياطي الطارئ المسموح. & \textenglish{public MotherShip(int fleetNumber, int baseNumber, int shipNumber, string captainID, string fleetName, double emergencyReserve)} \\
    \hline
    عمليات \textenglish{Get} للصفات الخاصة بها. & \textenglish{GetFleetName(), GetEmergencyReserve()} \\
    \hline
    عملية \textenglish{Set} لصفة اسم الأسطول، وعملية \textenglish{Set} لصفة حجم الاحتياطي الطارئ المسموح (يجب أن يكون غير سالب). & \textenglish{SetFleetName(string fleetName)} \newline \textenglish{SetEmergencyReserve(double reserve)} \\
    \hline
    عملية تصف الكائن. & \textenglish{ToString()} \\
    \hline
    عملية تعبئة وقود (موروثة من \textenglish{BasicShip}؛ الكمية يجب أن تكون موجبة). تعيد \textenglish{true} إذا نجحت التعبئة، وإلا \textenglish{false}. & \textenglish{Refuel(double amount)} \\
    \hline
    عملية استهلاك وقود إذا كان مسموحاً (الكمية موجبة ولا يهبط الوقود تحت \textenglish{-emergencyReserve})، وتعيد \textenglish{true} إذا نجحت و \textenglish{false} إذا لم تنجح. & \textenglish{Consume(double amount)} \\
    \hline
    عملية نقل وقود إلى مركبة استكشاف تتلقى كمية الوقود المراد نقلها وكائناً من النمط \textenglish{ScoutShip}. تنقل العملية كمية الوقود من السفينة الأم الحالية إلى مركبة الاستكشاف المُتلقّاة. تعيد العملية \textenglish{true} إذا نجح النقل، و \textenglish{false} إذا لم ينجح بسبب كون المركبة \textenglish{null} أو كمية غير موجبة أو كمية تتجاوز ما تسمح به السفينة الأم من استهلاك. & \textenglish{TransferFuel(double amount, ScoutShip target)} \\
    \hline
    \end{tabular}
    \par} %

    \vspace{2em}
    \noindent \textbf{الفحص في \textenglish{Main}:} \\
    أنشئوا سفينة أم، عبّئوها بالوقود، وجرّبوا نقل وقود منها إلى مركبة استكشاف عبر \textenglish{TransferFuel} في حالة ناجحة وأخرى فاشلة، واطبعوا حالة المركبتين.

\end{enumerate}

\vspace{1cm}
\begin{flushleft}
أرجو لكم وقتًا ممتعًا.

الأستاذ محمود اغبارية.
\end{flushleft}

\end{document}
